\section{Discussion}
\label{sec:discussion}

The primary question is narrow: when multi-turn tutoring causes
ADHD-supportive scaffolds to slip, does second-pass oversight improve
structural adherence over a structure policy in the prompt alone?

Study~1 gives a three-part answer. Unconstrained tutoring almost never produces
the required scaffolds (\textbf{0\%}). The ADHD Assist policy recovers most of
the structure ($\sim$\textbf{76\%} turn-aware overall; $\sim$\textbf{86\%}
late). Oversight adds a further, measured but modest lift
($\sim$\textbf{80\%}~/~$\sim$\textbf{89\%} late), with larger gains on residual
hard turns. The aggregate gain from oversight is modest; the more important
contribution is architectural.

The takeaway is not that oversight replaces prompting. Reliable
ADHD-supportive structure is a \textbf{control stack}: written policy,
turn-aware routing, and an optional auditor that can revise a draft before the
learner sees it. Prompting does most of the work. Oversight is residual
insurance for multi-turn failure modes the prompt does not fully close. For a
tutor that must keep ADHD-supportive structure on by default, that stack is
the contribution: steering generative replies at the interface, with model,
retrieval, tools, and decoding held fixed.

We measured scaffold presence under automated checks. Hallucination motivated
the desire for control but was not a primary dependent variable. The human
pilot ($n{=}6$; Table~\ref{tab:pilot}) shows the protocol is operable and
paired descriptives lean Assist-ward; that is feasibility evidence, not a
population claim. Whether Assist helps students with ADHD \emph{more} than
their peers is a question for the powered follow-on study.

The modest oversight lift still supports the enforcement thesis, provided the
claim is kept at its actual size. We do not claim near-perfect overall rates,
confirmatory load reduction, or ADHD exclusivity.

Model capacity is a separate follow-on question. Study~1 freezes one drafting
model and uses that same model for Dean rewrites. Exploratory Assist-on runs
outside the freeze (local Qwen~7B vs.~32B, alongside the Flash baseline;
partial matrix, not a second frozen suite) suggest stronger models follow the
full ADHD policy more reliably during first-pass generation. Capacity
therefore looks most consequential at the pedagogical generation stage.
Future work should test whether narrower oversight can be delegated to smaller
models without sacrificing compliance, establishing a quality--latency--cost
frontier, rather than assuming every role needs the largest available model.
